\begin{proposition}[crossing and multiple-fibre modules]
\label{prop:relative-local}
The symbols $R_2$ and $R_3$ denote copies of the same regular local ring,
with different maps from the one-dimensional base:
let $R_2=\C\{x,y,z\}$ with base coordinate $t=xy$, and let
$R_3=\C\{x,y,z\}$ with base coordinate $t=xyz$.  Then
\[
\begin{aligned}
 M_2&:=\Omega^1_{R_2/\C\{t\}}
 =\frac{R_2\,dx\oplus R_2\,dy\oplus R_2\,dz}
        {R_2\,(y\,dx+x\,dy)},\\[2pt]
 M_3&:=\Omega^1_{R_3/\C\{t\}}
 =\frac{R_3\,dx\oplus R_3\,dy\oplus R_3\,dz}
        {R_3\,(yz\,dx+xz\,dy+xy\,dz)}.
\end{aligned}
\tag{6.1}
\]
Both modules are torsion-free, but neither is locally free at the origin.
They embed in free modules by
\[
\begin{aligned}
 M_2&\hookrightarrow R_2e\oplus R_2dz,
 &dx&\longmapsto xe,&dy&\longmapsto-ye,&dz&\longmapsto dz,\\
 M_3&\hookrightarrow R_3e_1\oplus R_3e_2,
 &dx&\longmapsto xe_1,&dy&\longmapsto ye_2,&dz&\longmapsto-z(e_1+e_2),
\end{aligned}
\tag{6.2}
\]
with images
\[
 R_2dz\oplus (x,y)e,
 \qquad
 R_3xe_1+R_3ye_2+R_3z(e_1+e_2).
\tag{6.3}
\]
If $R_m=\C\{s,y_1,y_2\}$ and the base coordinate is $t=s^m$ with
$m\ge2$, then
\[
\Omega^1_{R_m/\C\{t\}}
\simeq R_m\,dy_1\oplus R_m\,dy_2
       \oplus R_m/(s^{m-1})\,ds,
\tag{6.4}
\]
and the final summand is exactly its torsion submodule.
\end{proposition}

\begin{proof}
Differentiating the three equations gives the presentations in (6.1).  For
$M_2$, the first map in (6.2) sends $(a,b,c)$ to $(ax-by,c)$.  If this is
zero, then $c=0$ and $ax=by$.  Since $x$ and $y$ are relatively prime in the
UFD $R_2$, there is an $h$ with $b=xh$ and $a=yh$; hence its kernel is exactly
the displayed relation.  For $M_3$, the second map sends $(a,b,c)$ to
$(ax-cz,by-cz)$.  If it is zero, relative primeness first gives $x\mid c$,
then $y\mid(c/x)$; writing $c=xyh$ gives
$(a,b,c)=h(yz,xz,xy)$, again exactly the defining relation.  Thus both maps
are injective, proving torsion-freeness and (6.3).

At the maximal ideal $(x,y,z)$ each defining relation has all coefficients in
that maximal ideal.  Consequently the quotient of either module by the
maximal ideal has dimension $3$ over $\C$, while its generic rank is $2$;
a locally free module of rank $2$ would need only $2$ generators.  This proves
non-local-freeness at the origin.  Finally, for $t=s^m$ one has
$dt=m s^{m-1}ds=0$ relatively.  The integer $m$ is a unit in $\C$, so the
decomposition (6.4) follows.  Its last summand is annihilated by
$s^{m-1}$, while the other two summands are free; after inverting $s$ it
vanishes.  It is therefore exactly the torsion submodule.
\end{proof}

\begin{remark}[scope of the local correction]
\label{rem:relative-local-scope}
The modules in (6.1) are modules on the smooth total-space chart, and their
torsion-freeness is asserted before restricting to a fibre.  They must be
distinguished both from the intrinsic sheaf $\Omega^1_S$ of the singular
fibre and from the restricted relative module
$M_i\otimes_{R_i}R_i/(t)$ ($i=2,3$): tensoring with the fibre ring can create
conductor-supported torsion, as the $\tau$ calculation in the proof of
Theorem~\ref{thm:nodal-test} shows.  If a local quotient is finite
surjective \'{e}tale over a chart in (6.4), relative differentials commute with
this base change and faithful flatness then reflects torsion of the pulled-back
module; a ramified quotient
requires a separate invariant-module calculation.  Thus the bare $t=xyz$
chart has a torsion-free but non-locally-free total-space relative module for
$p=1$; this fact alone does not decide torsion after fibre restriction or
quotient.  The multiple-fibre charts do have the total-space torsion in
(6.4).  No claim about torsion in $\bigwedge^2\Omega^1$ or
$\bigwedge^3\Omega^1$, and no global direct-image or Hodge conclusion, is
made here.
\end{remark}

\section{\texorpdfstring{Relation to the proposed $S^6$ reconstruction}{Relation to the proposed S6 reconstruction}}

The larger validate/repair/disprove workbench applies the filtration above to
a proposed compact complex threefold with a non-normal singular fibre and
records candidate values
\[
 H^1(X;\Z)=H^2(X;\Z)=H^3(X;\Z)=0,
 \qquad a(X)=1,
 \qquad \langle c_3(T_X),[X]\rangle=2.
 \tag{7.1}
\]
It also records candidate-specific calculations of integral specialization,
Leray transgressions, and the nonvanishing of
$R^2f_*(T_X\otimes L)$.  Those calculations are the source manuscript's
application and are not repeated here.  None of the upstream construction,
topology, or Picard assertions is proved in this short note.  Applying Proposition
\ref{prop:criterion} to that candidate would require, at minimum, a complete
independent verification of the smooth compact threefold, the full reduced
Cartier fibre, the normalization and line-bundle restrictions, and the
actual connecting map $\delta_A$ (or an equivalent explicit computation).
The present result supplies the local algebraic target for that check; it is
not an adjudication of the candidate's existence.

\section*{AI-use and authorship disclosure}

Generative-AI systems assisted the transcription of source material, symbolic
and exact-arithmetic checks, literature searching, and preparation of this
draft.  The local drafting and audit assistance used OpenAI Codex gpt-5.6-sol,
Ultra; no AI system is an author.  A human author must verify every statement,
secure the right to submit the work, and take responsibility for the final
version, consistent with the current arXiv policy for generative-AI tools
\cite{arXivModeration}.

\bibliographystyle{amsplain}
\bibliography{references}

\end{document}
